\bookchapter{제12장: 영상/음악 엔터테인먼트와 독서}

\chapterintro{스마트폰은 지금 기본적으로 MP3, TV 리모컨, Kindle, 휴대용 게임기를 대체했습니다. 솔직히, 많은 사람이 이런 App의 가장 기본적인 기능만 사용하고, 아직 발견하지 못한 좋은 점이 많습니다. 이 장에서는 스마트폰 엔터테인먼트를 더 순조롭고 편안하게 만드는 방법을 이야기하겠습니다.}

\booksection{일,음악 듣기: 어떤 App을 고를지, 어떻게 들어야 편한지}

음악 App은 지금 기능이 다 비슷하지만, 각각 강조점이 있어, 잘 고르면 신경을 많이 쓸 필요가 없습니다:

\begin{itemize}
  \item \textbf{네이시윈 음악}: 강점은 플레이리스트 커뮤니티와 댓글란입니다. "비 오는 날 듣기 좋은 노래"나 "공부 지칠 때 쉬기 위한 순음악"을 검색하면 다른 사람이 정리한 플레이리스트가 있습니다. 댓글란이 때로는 노래보다 더 정감 있습니다. 단점은 일부 곡의 저작권이 불완전하고, 어떤 곡은 듣다 보면 회색으로 변합니다.
  \item \textbf{QQ음악}: 저작권이 가장 완전하고, 주걸륜, 오월일 같은 인기 가수의 곡은 기본적으로 QQ음악에만 있습니다. 음질 옵션도 풍부하고, 회원은 무손실 음질을 켤 수 있습니다. "듣고 싶은 노래를 모두 검색할 수 있는지"를 중요하게 생각한다면, QQ음악이 가장 안정적인 선택입니다.
  \item \textbf{Apple Music}: 애플 사용자의 첫 선택. 인터페이스가 깔끔하고, 광고가 없으며, 음질이 좋고, 무손실과 Dolby Atmos를 지원합니다. 하지만 추천 알고리즘은 국내 App만큼 중국어 사용자의 취향을 잘 이해하지 못하고, 플레이리스트 문화도 네이시윈처럼 발달하지 않았습니다. 간결한 경험을 좋아하고, 소셜 기능을 크게 신경 쓰지 않는 사람에게 적합합니다.
  \item \textbf{Spotify}: 글로벌 사용자가 가장 많고, 추천 알고리즘이 공인된 최강입니다. "주간 발견"과 "일일 추천"의 정확도가 놀라울 정도로 높아, 오래 쓸수록 당신을 더 잘 이해합니다. 단점은 무료 버전은 광고가 많고, 일부 중국어 곡의 저작권 커버가 불완전합니다. 사용하려면 특수 네트워크가 필요합니다.
\end{itemize}

\textbf{몇 가지 음악 감상 소팁}: 만약 Bluetooth 이어폰을 사용한다면, 음악 App 설정에서 “음향 효과 조절”이나 “이퀄라이저”가 있는지 확인하고, 저음과 고음을 살짝 조정하면 효과가 크게 향상됩니다. 또한, 스트리밍 음악 감상은 데이터를 많이 소모하므로, 설정에서 “WiFi 환경에서만 무손실 재생”을 켜거나, 미리 좋아하는 플레이리스트를 다운로드하여 오프라인으로 들으세요.

\booksection{이,팟캐스트(Podcast): 출퇴근 길의 “귀 동반자”}

팟캐스트라는 것은, 들어보지 않은 사람은 자신과 무관하다고 생각하지만, 한번 듣기 시작하면 돌아갈 수 없습니다. 본질은 “라디오 프로그램의 Netflix판”—과학기술을 이야기하는 사람, 역사를 이야기하는 사람, 이야기를 읽어주는 사람이 있어, 어떤 유형이든 다 있습니다.

\begin{itemize}
  \item \textbf{샤오위저우 App}: 국내에서 가장 잘 만들어진 팟캐스트 App, 인터페이스가 깔끔하고, 잡다한 기능이 없습니다. "매일 추천"과 "편집 정선"이 새로운 프로그램을 발견하게 해주고, 댓글란도 활발합니다. 제가 가장 많이 듣는 것은 "좌충우돌"(문화 역사류)과 "부합시의"(사회 이슈류)이며, 매 회 똑똑한 사람들이 이야기하는 것을 방청하는 것 같습니다.
  \item \textbf{애플 팟캐스트}: iPhone에 내장되어, 추가 App 다운로드가 필요 없습니다. CarPlay와의 통합이 가장 좋아, 운전할 때 편리합니다.
  \item \textbf{청취 기술}: 대부분의 팟캐스트 App은 1.2배속 또는 1.5배속 재생을 지원하며, 익숙해지면 적지 않은 시간을 절약할 수 있습니다. 저는 일반적으로 중국어는 1.2배, 영어는 1.5배로 듣습니다. 또한 반드시 "새 프로그램 자동 다운로드"와 "재생 후 자동 삭제"를 켜서, 스마트폰에 들은 오디오 파일이 쌓이지 않도록 하세요.
\end{itemize}

\booksection{삼,영상 보기: 재생기에서 스트리밍까지 올바른 열기 방식}

많은 사람이 모르는 사실: 스마트폰에서 재생할 수 있는 형식은 스트리밍 App에서 제공하는 것보다 훨씬 많습니다. 좋은 로컬 재생기가 많은 문제를 해결할 수 있습니다.

\begin{itemize}
  \item \textbf{VLC 재생기}: 오픈소스 무료, 거의 모든 형식을 재생—MKV, AVI, FLAC, 심지어 ISO 디스크 이미지. 만약 다운로드한 영상이 컴퓨터에서는 재생되지만 스마트폰에서는 재생되지 않는다면, 십중팔구 형식 문제이며, VLC를 설치하면 해결됩니다. 외장 자막도 지원하여, .srt 자막 파일과 영상을 같은 폴더에 두면 VLC가 자동으로 로드합니다.
  \item \textbf{nPlayer}: 더 정교한 경험을 원한다면, 이 유료 App은 살 만합니다. 인터페이스가 보기 좋고, 제스처 컨트롤이 매우 편리—왼쪽 상하 슬라이드로 밝기 조절, 오른쪽 상하 슬라이드로 볼륨 조절, 좌우 슬라이드로 빨리 감기/되감기. DTS와 Dolby 사운드 효과를 지원하여, 대작 시청 시 현장감이 훨씬 강합니다.
  \item \textbf{Bilibili(B스테이션)}: 더 이상 2차원 사이트가 아니며, 지금 학습 자원이 엄청나게 많습니다. 대학원 입시, 프로그래밍, 사진, 요리, 어떤 튜토리얼이든 다 있습니다. 탄막 문화는 B스테이션의 영혼이며, 영상 볼 때 탄막을 켜면, 여러 사람이 함께 보는 느낌입니다. 다만 학습 영상 시청 시에는 탄막을 끌 것을 권장하며, 그렇지 않으면 집중하기 어렵습니다.
  \item \textbf{TV 화면 미러링}: 스마트폰 화면이 아무리 커도 TV만큼 편하지 않습니다. 같은 WiFi에서, 영상 App 오른쪽 상단의 "TV" 화면 미러링 버튼을 누르면, 화면을 TV에 투사할 수 있습니다. 대부분의 스마트 TV와 프로젝터가 지원하며, 드라마나 영화 시청 경험이 훨씬 좋습니다. iPhone 사용자는 AirPlay로 전체 화면을 직접 미러링할 수도 있습니다.
\end{itemize}

\booksection{사,전자책 독서: 스마트폰에서도 잘 책 읽을 수 있습니다}

누군가는 스마트폰으로 책을 읽으면 눈이 상하고, 집중하기 어렵다고 합니다, 이는 사실 일리가 있습니다. 하지만 스마트폰 독서의 장점도 분명합니다—언제 어디서나, 파편화된 시간을 최대한 활용. 저는 지하철에서, 줄 서 있을 때, 자기 전에 스마트폰으로 적지 않은 책을 읽었습니다.

\begin{itemize}
  \item \textbf{위챗 독서}: 현재 국내에서 경험이 가장 좋은 독서 App. 서고가 매우 완전하며, 출판물과 웹소설 모두 있습니다. "독서 시간으로 책 코인 교환" 모델을 사용하여, 많이 읽으면 무료로 책을 교환할 수 있어 학생에게 친화적입니다. "밑줄" 기능으로 다른 사람이 어떤 단락에 표시했는지 볼 수 있으며, 때로는 다른 사람이 표시한 좋은 문장을 보고 공감하게 됩니다. 다만 "독서 순위"는 좀 지나쳐서, "시간을 채우기 위해 책을 읽는" 것이 되기 쉬워, 끄기를 권장합니다.
  \item \textbf{Kindle App}: 만약 Kindle 기기가 있다면, 스마트폰용 Kindle App으로 읽기 진행도를 동기화할 수 있습니다. Kindle의 정품 전자책 레이아웃 품질은 해적판 PDF보다 훨씬 좋지만, Kindle 중국 서점은 이미 닫혔고, 해외 지역으로 전환해야 합니다.
  \item \textbf{Apple Books(도서)}: iPhone에 내장되어, EPUB과 PDF 읽기에 적합. 인터페이스가 극도로 간결하고, 페이지 넘기기 애니메이션이 편안합니다. 인터넷에서 전자책 파일을 다운로드한 경우, "도서로 공유"로 바로 가져오기할 수 있습니다.
  \item \textbf{독서 설정 권장}: 어떤 App을 사용하든, 반드시 배경을 따뜻한 노란 톤으로 설정하고(눈 보호 모드), 폰트를 약간 키우고, 행간을 넓히세요. 밤에 책 읽을 때는 "다크 모드"를 켜세요—검은 바탕에 흰 글자, 눈이 훨씬 편합니다. 절대 어두운 빛 아래에서 흰 바탕에 검은 글자, 작은 폰트로 오래 읽지 마세요, 눈 피로는 누적됩니다.
\end{itemize}

\booksection{오,스마트폰 게임: 즐겁게즐기기, 스마트폰도 상하지 않게}

저는 모바일 게임에 대해 다음과 같은 입장입니다: 놀아도 좋지만, 게임이 당신을 통제하게 두지 마세요. 게임 도구를 합리적으로 사용하면, 더 즐겁게 놀 수 있고, 정당한 일도 지장받지 않습니다.\textbf{게임 모드/게임 공간}: 현재 주류 스마트폰에는 이 기능이 있습니다. 켜면 스마트폰이 자동으로 백그라운드를 정리하고, 알림을 차단하며, 오터치 방지—팀전 중인데 갑자기 전화가 뜨는 일을 게임 모드가 차단해 줍니다. 일부 스마트폰은 성능 모드/E-sports 모드도 지원하여, 프로세서를 강제로 최대로 돌리지만, 대가는 발열과 배터리 소모가 더 빨라집니다.\textbf{화질과 프레임률 어떻게 조정}: 무조건 최대로 올리지 마세요. 먼저 프레임률(부드러움)을 보장하고, 그 다음 화질을 추구하세요. 60프레임은 30프레임보다 훨씬 편하지만, 120프레임은 스마트폰에 큰 부담이 되어, 중급기는 강요할 필요 없습니다. 만약 경쟁류 게임(예: 왕자영요, 배틀그라운드)을 한다면, 프레임률 우선순위가 화질보다 훨씬 높습니다—부드러운 조작감이 좋은 잔디보다 100배 중요합니다.\textbf{게임 중 스마트폰이 뜨거워지면 어떡해}: 이는 정상이지만, 장기간 고온은 배터리 수명을 손상시킵니다. 몇 가지 온도 낮추는 방법: 스마트폰 케이스를 벗겨 열 방출; 화면 밝기를 70\% 정도로 낮춤(밝기는 발열원); 충전하면서 게임하지 마세요(발열 가중, 배터리 빨리 사망). 만약 스마트폰이 손으로 잡을 수 없을 정도로 뜨거워지면, 즉시 종료하고 잠시 쉬게 하세요.\textbf{과몰입 방지 소도구}: iOS의 “화면 사용 시간”과 안드로이드의 “디지털 건강”은 게임 App에 매일 시간 한도를 설정할 수 있습니다. 제가 설정한 것은 평일 최대 1시간, 주말 2시간입니다. 시간이 되면 App 아이콘이 회색으로 변하고, 클릭해도 들어가지 않습니다. 이 제한은 정말 유용하며, 그저 형식적인 것이 아닙니다.\textbf{클라우드 게임}: 만약 스마트폰 성능이 부족하지만 대작을 플레이하고 싶다면, 클라우드 게임 플랫폼(예: 텐센트 START 클라우드 게임, 네테이즈 클라우드 게임)을 시도해 보세요. 게임은 원격 서버에서 실행되고, 화면은 영상 스트림 방식으로 스마트폰에 전송됩니다. 네트워크 요구사항은 높지만 스마트폰 성능은 거의 요구하지 않습니다. 단점은 지연감을 피할 수 없어, 전략류, RPG류에 적합하고, FPS 사격류에는 부적합합니다.직접 해보기: 당신의 스마트폰 엔터테인먼트 시스템 구축하기음악 App에서 가장 좋아하는 3개 앨범이나 플레이리스트를 다운로드하여, 오프라인에서 들을 수 있게 하세요.샤오위저우나 애플 팟캐스트에서 관심 있는 3개 채널을 구독하고, 한 에피소드를 다운로드하여 시청해 보세요.VLC 재생기를 설치하고, 컴퓨터에서 복사해 온 영상 파일을 재생해 보세요.위챗 독서에서 읽고 싶은 책을 찾고, 따뜻한 배경과 눈 보호 모드를 설정하며, 폰트를 자신에게 맞는 크기로 조정하세요.만약 게임을 한다면, 게임 모드를 켜서 어떤 설정을 최적화할 수 있는지(방해 금지, 프레임률 우선 등) 확인해 보세요.
